\section{Methodology}
\label{sec:method}

Figure~\ref{fig:pipeline} gives a schematic overview of the integrated framework. A sliding window of 30 time steps over the 96 sensor channels is the common input representation. Stage 1 produces a binary anomaly flag, stage 2 produces an 18-class accident label plus per-sensor and per-timestep attributions, and stage 3 produces a ten-step forward trajectory prediction along with the scalar physics residuals attached to that prediction. The downstream consistency check combines the stage-2 label with the stage-3 trajectory, as detailed in Section~\ref{sec:integration}.

\begin{figure}[H]
    \centering
    \placeholder{Figure 1: pipeline schematic. Boxes for modules 1--3 with arrows showing data flow; inset showing the cross-validation check between module 2 and module 3.}
    \caption{Overview of the three-stage integrated framework. Stage 1 performs semi-supervised screening on normal-data-trained models. Stage 2 performs supervised 18-class classification with SHAP attribution. Stage 3 produces physics-informed short-horizon trajectory predictions that can be cross-checked against the stage-2 classification.}
    \label{fig:pipeline}
\end{figure}

\subsection{Data Representation}
\label{sec:data-rep}

Raw \nppad{} recordings are resampled onto a uniform time grid and concatenated into per-scenario sequences of $96$ sensor channels covering primary-loop pressures, hot/cold-leg temperatures, coolant flow rates, reactor power, neutron flux, pressurizer level, steam generator levels and pressures, and boron concentration. Sequences are split by scenario into $70/15/15$ train/validation/test partitions to avoid leakage within a single transient. A standard scaler is fit on the normal-operation training subset only, and applied to all partitions.

Sliding windows of length $T = 30$ time steps are drawn with stride $s_{\text{train}} = 10$ during training and $s_{\text{eval}} = 1$ during evaluation, yielding $32{,}284$ training windows and $7{,}316$ test windows. Each window is labelled with both a binary anomaly flag and the 18-way accident-type label of its parent scenario.

\subsection{Stage 1: Semi-Supervised Anomaly Detection}
\label{sec:anomaly}

Stage 1 comprises three architectures trained only on the 28 normal-operation windows in the training set. Anomaly scores are derived from reconstruction or prediction error, and a decision threshold is tuned on validation data to maximize F1.

\paragraph{Autoencoder.} A fully connected encoder--decoder with layer widths $96 \rightarrow 64 \rightarrow 32 \rightarrow 16 \rightarrow 32 \rightarrow 64 \rightarrow 96$, ReLU activations, and batch normalization reconstructs each timestep independently. The anomaly score for a window is the mean per-timestep mean-squared reconstruction error.

\paragraph{LSTM next-step predictor.} A bidirectional LSTM with hidden size 128 is trained with teacher forcing to predict timestep $t+1$ from timesteps $1{:}t$. The anomaly score is the mean squared prediction error across the window.

\paragraph{Masked Transformer reconstructor.} A Transformer encoder with sinusoidal positional encoding, three layers, and eight attention heads reconstructs a randomly masked $15\%$ of the timesteps in each training window. At inference, the model is asked to reconstruct every timestep, and the mean reconstruction error is used as the anomaly score.

All three models are trained with Adam, a learning rate of $10^{-3}$, early stopping on validation loss with patience 10, and a batch size of 64.

\subsection{Stage 2: Accident Classification with Root-Cause Attribution}
\label{sec:classifier}

Stage 2 is a supervised 18-class classifier operating on the same 30-step window representation. We train three architectures to assess how inductive bias affects performance on this imbalanced problem.

\paragraph{1D-CNN classifier.} Three convolutional blocks ($96\rightarrow 64 \rightarrow 128 \rightarrow 256$ channels, kernel size 3, ReLU, batch normalization) feed into a global average pooling layer, a 128-unit fully connected layer, and an 18-way softmax head. Dropout of $0.3$ is applied before the final layer.

\paragraph{Bidirectional LSTM classifier.} A two-layer bidirectional LSTM with hidden size 128 is unrolled over the 30-step window. The final hidden states from both directions are concatenated ($256$ units), passed through a 128-unit fully connected layer with dropout, and then into the softmax head.

\paragraph{Transformer classifier.} A learnable \texttt{[CLS]} token is prepended to the window, positional encodings are added, and three encoder layers (model dimension 128, eight heads, feed-forward dimension 512) attend over the sequence. The final \texttt{[CLS]} representation is passed through LayerNorm, a 128-unit linear layer with ReLU and dropout, and the softmax head.

All classifiers are trained with cross-entropy loss, Adam, early stopping with patience 15, and learning rates $10^{-3}$ for CNN/LSTM and $5 \times 10^{-4}$ for the Transformer.

\paragraph{Root-cause attribution.} We compute two complementary attribution signals on the best-performing classifier. SHAP values~\citep{lundberg2017shap} are computed with the \texttt{DeepExplainer} backend using a 100-window background sample drawn from the training set. This yields a per-sample tensor of shape $(T \times F)$ assigning an importance score to every (timestep, sensor) pair. We additionally compute input-$\times$-gradient attributions as a faster alternative for large-scale analyses. Per-class feature importance is obtained by averaging absolute SHAP values over all test samples of that class; per-class temporal importance is obtained by averaging over sensors. In Section~\ref{sec:results-shap} we map the top-ranked sensors of representative classes onto physical parameter groups (primary pressure, coolant flow, steam generator secondary side, \textit{etc.}) and compare against the known mechanistic signature of each accident.

\subsection{Stage 3: Physics-Informed Trajectory Prediction}
\label{sec:pinn}

Stage 3 is a GRU encoder--decoder surrogate trained under a composite physics-informed loss. Given a context window of $T_c = 20$ time steps, the model predicts the next $T_p = 10$ time steps of all 96 sensor channels.

\subsubsection{Architecture}

The encoder is a two-layer GRU with hidden size 128 that ingests the context window and produces a final hidden state $\mathbf{h}_{T_c}$. The decoder is an autoregressive GRU initialized from $\mathbf{h}_{T_c}$ that emits, at each decode step $k = 1, \dots, T_p$, a delta vector $\Delta \hat{\mathbf{x}}_k \in \mathbb{R}^{96}$ via a linear projection. The predicted state is obtained by a residual connection
\begin{equation}
    \hat{\mathbf{x}}_{T_c + k} = \mathbf{x}_{T_c} + \sum_{j = 1}^{k} \Delta \hat{\mathbf{x}}_j,
    \label{eq:residual}
\end{equation}
which biases the model towards the local linearization of the dynamics and stabilizes autoregressive roll-outs. The total parameter count is $384$K, small enough for real-time inference on a CPU.

\subsubsection{Physics-Informed Loss}

The training loss combines four terms:
\begin{equation}
    \mathcal{L}(\theta) = \lambda_{\text{data}} \mathcal{L}_{\text{data}} + w(t)\bigl[\lambda_{\text{pk}} \mathcal{L}_{\text{pk}} + \lambda_{\text{cons}} \mathcal{L}_{\text{cons}}\bigr] + \lambda_{\text{bnd}} \mathcal{L}_{\text{bnd}},
    \label{eq:loss}
\end{equation}
where $w(t)$ is a linear warmup schedule ramping from $0$ to $1$ over the first $500$ training steps. We set $\lambda_{\text{data}} = 1.0$, $\lambda_{\text{pk}} = 0.1$, $\lambda_{\text{cons}} = 0.05$, $\lambda_{\text{bnd}} = 0.01$.

\paragraph{Data term.} $\mathcal{L}_{\text{data}}$ is the mean-squared error between predicted and true sensor trajectories in the standard-scaler-normalized space.

\paragraph{Point-kinetics residual.} Let $n(t)$ denote the predicted neutron density channel and let $T_{\text{avg}}(t)$ and $C_{\text{B}}(t)$ denote the predicted coolant average temperature and boron concentration. A simplified point kinetics model assumes
\begin{equation}
    \frac{dn}{dt} = \frac{\rho(t)}{\Lambda} n(t),
    \qquad
    \rho(t) = \alpha_T (T_{\text{avg}}(t) - T_{\text{ref}}) + \alpha_B (C_{\text{B}}(t) - C_{\text{B,ref}}),
    \label{eq:point-kinetics}
\end{equation}
with prompt neutron generation time $\Lambda = 2.0 \times 10^{-5}$~s, temperature coefficient $\alpha_T = -2.0 \times 10^{-4}$~$\Delta k/k$/\textdegree{}C, boron coefficient $\alpha_B = -1.0 \times 10^{-4}$~$\Delta k/k$/ppm, reference temperature $T_{\text{ref}} = 310$~\textdegree{}C, and reference boron concentration $C_{\text{B,ref}} = 200$~ppm, consistent with standard large-PWR parameters~\citep{duderstadt1976nuclear}. The residual is the squared error between the finite-difference time derivative of the predicted $n(t)$ and the right-hand side of Eq.~\eqref{eq:point-kinetics} evaluated on the prediction, averaged over the prediction horizon and normalized by a characteristic neutron density scale.

\paragraph{Energy-conservation residual.} At steady state, core thermal power $Q$ is related to primary loop flow $W$, coolant specific heat $c_p = 5.5$~kJ/kg/\textdegree{}C, and the hot/cold leg temperature difference $\Delta T = T_{\text{hot}} - T_{\text{cold}}$ via
\begin{equation}
    Q \approx W \, c_p \, \Delta T.
    \label{eq:conservation}
\end{equation}
Using the predicted values of $Q$, $W$, $T_{\text{hot}}$, and $T_{\text{cold}}$, $\mathcal{L}_{\text{cons}}$ is the squared residual of Eq.~\eqref{eq:conservation} normalized by a characteristic power scale. During rapid transients Eq.~\eqref{eq:conservation} does not hold exactly, but its inclusion as a soft penalty biases the surrogate towards quasi-steady energy balance, which is physically correct on sufficiently long time scales.

\paragraph{Physical-bounds penalty.} For each channel in a predefined set (pressures, power, temperatures, flow rates, boron concentration), $\mathcal{L}_{\text{bnd}}$ applies a one-sided quadratic penalty whenever the prediction falls below the physically admissible floor.

\subsubsection{Solver}

Where the point-kinetics residual is integrated rather than differentiated during evaluation, we use a fully differentiable fourth-order Runge--Kutta integrator implemented in PyTorch, enabling gradient flow through physical rollouts if desired for future extensions.

\subsection{Pipeline Integration and Cross-Validation}
\label{sec:integration}

At deployment, a new sliding window enters the pipeline and is first passed through stage 1. If stage 1 flags the window as anomalous, it is forwarded to stage 2 for 18-class classification and to stage 3 for trajectory prediction. Stage 2 also returns SHAP attributions for the predicted class. Stage 3 returns the 10-step forecast together with the associated point-kinetics and conservation residuals evaluated on the prediction.

We define a \emph{physics-consistency score} $\kappa$ as a weighted sum of the point-kinetics residual and the conservation violation on the predicted trajectory, normalized so that $\kappa \approx 1$ corresponds to the median residual observed on correctly classified test windows and larger values correspond to worse physical consistency. A threshold $\kappa^\star$ is calibrated on the validation set; windows for which $\kappa > \kappa^\star$ are flagged as low-confidence classifications and escalated for human review.

This mechanism is conservative: it does not attempt to correct the classifier, but it provides an auditable, physics-grounded signal that can reduce the rate of silent misclassification. Section~\ref{sec:results-pipeline} evaluates the effectiveness of this consistency check empirically.
